% TOP_PROBLEM: {"id": 3353, "title": "Polignac's Conjecture", "category_id": 1, "category": {"id": 1, "name": "number_theory", "display_name": "Number Theory", "description": "Properties of integers, prime numbers, Diophantine equations.", "slug": "number-theory", "order_index": 1, "created_at": "2026-07-31T15:26:25.670Z"}, "status": "open", "problem_number": "OPG-37289", "set_id": 12}
% FIRST_POSTED: 2026-10-01
% ENTRY_KIND: statement_only
% UnsolvedMath ID 3353; source catalogue code OPG-37289
% !TeX program = lualatex
% Definition and sources only; this file is not an LLM solution attempt.
\documentclass[11pt,a4paper]{article}
\usepackage[margin=25mm]{geometry}
\usepackage{fontspec}
\setmainfont{lmroman10-regular.otf}[BoldFont=lmroman10-bold.otf,
 ItalicFont=lmroman10-italic.otf,BoldItalicFont=lmroman10-bolditalic.otf]
\usepackage{amsmath,amssymb,mathtools}
\usepackage{unicode-math}
\setmathfont{latinmodern-math.otf}
\AtBeginDocument{\let\setminus\smallsetminus}
\usepackage{xurl,hyperref}
\hypersetup{colorlinks=true,urlcolor=blue}
\setcounter{secnumdepth}{0}
\setlength{\parindent}{0pt}
\setlength{\parskip}{5pt}
\setlength{\emergencystretch}{3em}
\begin{document}
\section*{Polignac's Conjecture}\label{3353}
{\small UnsolvedMath ID 3353 · OPG-37289 · Number Theory · OpenGarden}
\subsection{Definitions and mathematical statement}
Let $p_1=2<p_2=3<p_3<\cdots$ be all primes in increasing order, and
write $g_j=p_{j+1}-p_j$ for the gap between consecutive primes.
For a fixed positive even integer $h$, let
\[
 P_h(x)=\#\{j\ge1:p_j\le x,\ g_j=h\}.
\]
Polignac's conjecture is
\[
 \forall h\in2\mathbb Z_{>0},\qquad
              \lim_{x\to\infty}P_h(x)=\infty.
\]
An equivalent quantified version says that for every such $h$ and
every $X$ there are primes $p>X$ and $p+h$ with no prime strictly
between them. The interval condition is essential: merely finding
infinitely many prime pairs of difference $h$ gives a weaker assertion
when $h>2$. The evenness restriction comes from the fact that all
primes other than two are odd.

Each gap $h$ is fixed in its own infinitude assertion; the bound beyond
which a pair is sought may be arbitrary. The conjecture does not claim
uniform estimates in $h$, nor a formula for $P_h(x)$. The case $h=2$
is the twin-prime conjecture. A disproof of the universal assertion
would identify an even $h$ occurring only finitely often as a consecutive
prime gap.
\subsection{Short English statement}
Does every positive even integer occur infinitely often as the gap between consecutive primes?
\subsection{Sources}
\begin{itemize}
\item[S1] ULAM.AI and the UnsolvedMath Contributors, supplied problems.json, record 3353 (OPG-37289), OpenGarden collection. Catalogue identity and source transcription; CC BY 4.0.
\url{https://huggingface.co/datasets/ulamai/UnsolvedMath}.
\item[S2] Open Problem Garden, Polignac's Conjecture. Original problem and discussion; the mathematical formulation below is expanded from the supplied source record.
\url{http://www.openproblemgarden.org/op/twin_primes_and_polignacs_conjecture}.
\end{itemize}
\end{document}
